\section{The formula firewall and the four-part interface}
\label{sec:firewall}

\subsection{Why an interface needs four parts}

A reusable component is often described by only a program and a semantic theorem.
That is sufficient in an ordinary programming library but insufficient for an
embedded weak theory.  We require four distinct fields.

\begin{definition}[FRNF component]
An \(\FRNF\) component from input presentation \(X\) to output presentation \(Y\)
consists of:
\begin{enumerate}[label=(\roman*),leftmargin=2.4em]
\item fixed primitive-recursive symbols computing finite output observations;
\item an object-language formula \(C_{X\to Y}\) expressing the component contract;
\item a constructor that, from embedded derivations of declared input contracts,
returns an embedded derivation of \(C_{X\to Y}\); and
\item an optional host-semantic theorem identifying the numerical contract with its
intended mathematical meaning.
\end{enumerate}
The first three fields are the \emph{kernel-facing interface}.  The fourth is the
\emph{semantic interpretation}.
\end{definition}

The separation between (iii) and (iv) is the formula firewall.  It may feel redundant
when the host theorem is elementary, but it allows components to be combined without
accidentally importing a strong metatheory into the object proof.

\subsection{Exact contracts}

The simplest component computes a term \(w(x)\) and proves a formula
\(Q(x,w(x))\) from a precondition \(P(x)\).  At the object level it has the shape
\[
    [P]\ \vdash_{\PRAAx}\ Q[w].
\]
Existential hiding then gives
\[
    [P]\ \vdash_{\PRAAx}\ \exists y\,Q(x,y),
\]
and implication introduction gives a closed contract.  Sequential composition is
substitution followed by cut.  The accompanying
\code{lean/FRNF/Kernel.lean} records exactly these logical patterns through a proof
schedule.

The crucial point is that an exact contract is not merely a structure containing a
Lean function \(w\) and a field of type \(Q(w)\).  Its witness is an object term
\(\PTm\), and its certificate is an embedded derivation.

\subsection{Relational contracts}

Many constructive constructions do not justify a single exact witness.  A polynomial
with a multiple root, an eigenvalue cluster, a chart overlap, or a discretized
optimization problem naturally yields a nonempty finite set of acceptable answers.
For such a component we use a bounded relation
\[
  R(x,i), \qquad i<B(x),
\]
with an object derivation of
\[
  P(x)\ \longrightarrow\ \exists i<B(x)\,R(x,i).
\]
The relation may carry a compatibility edge to the candidates at the next precision.
A downstream component must either be defined relationally on every valid candidate or
prove that its observable is stable across the candidate cloud.

This is stronger than choosing the least candidate immediately.  Least choice is
primitive recursive over a finite set, but it can destroy extensionality: two
presentations of the same represented object may enumerate candidates in different
orders.  Keeping the relation until the observation point allows one to prove
presentation-independent results.

\subsection{Conditional contracts and boundary nodes}

Suppose a finite algorithm enumerates all solutions below a supplied height bound
\(H(p)\).  Its correctness is a theorem of weak arithmetic once \(H\) is fixed and the
height theorem is given.  The statement should therefore be represented as
\[
  \operatorname{HeightBound}(p,H(p))
  \longrightarrow
  \operatorname{CompleteList}(p,\operatorname{enum}(p,H(p))).
\]
The height theorem is a boundary node.  It is not hidden in the implementation of
\code{enum}.  If a later project supplies an object proof of the height bound, ordinary
cut discharges the boundary.  Until then the endpoint remains a precise conditional
theorem rather than an incorrect claim of unconditional enumeration.

\begin{principle}[Producer/verifier separation]
A primitive-recursive verifier for a finite certificate does not become a
primitive-recursive producer of that certificate.  Production and verification are
separate nodes with separate proof obligations.
\end{principle}

This distinction is decisive in the general Siegel and Faltings packets, exceptional-
zero arguments in analytic number theory, arbitrary compactness arguments, and exact
splitting data for represented algebraic functions.

\subsection{Status fields as part of the mathematics}

A practical artifact should store, for every component, at least the following status
fields:

\begin{enumerate}[label=\arabic*.,leftmargin=2.2em]
\item semantic statement written;
\item finite reference implementation tested;
\item object formula defined;
\item local \(\Prf\,\PRAAx\) constructor implemented;
\item composed endpoint built;
\item native transfer built;
\item remaining boundary assumptions enumerated;
\item axiom audit inspected.
\end{enumerate}

These fields are not project-management decoration.  They express logically distinct
claims.  In particular, a compiled host theorem with a remaining \code{sorry} in the
object constructor is not a partial object proof; it is a complete result at one level
and an open result at another.

\subsection{Compositionality}

Let \(C_1,\ldots,C_m\) be component contracts arranged in topological order.  The local
proof for node \(i\) may use only the permanent hypotheses \(\Gamma\) and a declared
subset of earlier contracts.  If
\[
   \Delta_i,\Gamma\vdash_{\PRAAx} C_i
   \quad\text{and}\quad
   \Gamma\vdash_{\PRAAx}D\quad(D\in\Delta_i),
\]
then simultaneous cut yields \(\Gamma\vdash_{\PRAAx}C_i\).  Induction on the node list
therefore produces all graph contracts.  This simple observation is the logical core
of the proof schedule in \cref{sec:proof-schedule}; all substantive work lies in
making each local contract finite and QF-internalizable.
